\documentclass[addpoints]{exam}
% \documentclass[addpoints,12pt]{exam}

\usepackage[slovene]{babel}
\usepackage{amsfonts}
\usepackage{amssymb}
\usepackage{amsmath}
\usepackage{float}
\usepackage{graphicx}
\usepackage{enumitem}
\usepackage{color}
\usepackage{eurosym}

\usepackage{pgfpages}
% \usepackage{tikz}
\usepackage{wrapfig}
\usepackage{pgfkeys}
\usepackage{pgfplots}
\usepackage{xcolor}
\usepackage{tkz-euclide}
\usepackage{xfp}




%Komentar naslednje vrstice glede na verzijo testa -- rešitve ali prostor za odgovore:
% \printanswers

% \linespread{2}


\pointsinrightmargin
\marginbonuspointname{}


\renewcommand{\solutiontitle}{\noindent\textbf{Rešitve in točkovanje:}\par\noindent}

\colorgrids
\definecolor{GridColor}{gray}{0.7}
\setlength{\gridsize}{7mm}

\extrawidth{5mm}
\setlength{\rightpointsmargin}{1.5cm}

\begin{document}

\runningfooter{}{}{}

\begin{center}
    \fbox{\fbox{\parbox{15cm}{\centering
    Drugo pisno preverjanje znanja \\ 2. letnik \\ 12. december 2025 \\ (Kotne funkcije; Vektorji)}}}
\end{center}

\vspace{0.1in}
\makebox[\textwidth]{Ime in priimek, oddelek:\enspace\hrulefill}


\begin{center}
    \chqword{Naloga}
    \chpword{Možne točke}
    \chbpword{Dodatne točke}
    \chsword{Dosežene točke}
    \chtword{Skupaj}  
    \combinedpointtable[h][questions]
    % \combinedgradetable[h][questions]

\end{center}

\vspace{0.1in}


\begin{questions}

    % 1. vprašanje

    \question[3]
    Poenostavite izraz. 
    $$ \dfrac{\tan x}{\cos x - \frac{1}{\cos x}}+\sin x + \sin x \cot^2 x$$

    % \begin{solution}[6.5cm]
       
    % \end{solution}


    % 2. vprašanje

    \question[4]
    Izračunajte natančno vrednost izraza. 
    $$ \dfrac{\sin{120^\circ} + \cos{60^\circ}}{\sin{45^\circ} - \tan{300^\circ}}+\dfrac{\cos^2{310^\circ}+\sin^2{50^\circ}}{\sin{270^\circ}}$$

    % \begin{solution}[6cm]
       
    % \end{solution}


% \newpage
    % 3. vprašanje

    \question
    V kvadru $ABCDEFGH$ je točka $M$ razpolovišče roba $CD$, točka $O$ razdeli rob $EF$ v razmerju \\ $|EO|:|OF|=1:6$.
    Z baznimi vektorji $\overrightarrow{a}=\overrightarrow{AB}, \overrightarrow{b}=\overrightarrow{AD}$ in $\overrightarrow{c}=\overrightarrow{AE}$ izrazite vektorje: 
    \begin{parts}
        \part[2]
        $\overrightarrow{MO}$, 

    % \begin{solution}[5cm]

    % \end{solution}
        
        \part[3] 
        $\overrightarrow{BH}-\overrightarrow{GH}+\overrightarrow{GA}+\overrightarrow{CD}$.
    \end{parts}

    % \begin{solution}[5cm]

    % \end{solution}

    

    
    % 4. vprašanje

    \question[4]
    V trikotniku $\triangle ABC$ z oglišči $A(2,5,-3)$, $B(5,7,3)$ in $C(10,0,0)$ izračunajte koordinate težišča $T$ ter razdaljo med težiščem $T$ in ogliščem $A$.
    % \begin{parts}

    %     \part[2]
    %     dolžino težiščnice na stranico $AC=b$, 

    % \begin{solution}[6cm]

    % \end{solution}

    %     \part[3] 
    %     koordinate težišča trikotnika $\triangle ABC$.

    % \begin{solution}[2cm]

    % \end{solution}        
    
    % \part[5]
    %     skalarni produkt $\overrightarrow{AB}\cdot\overrightarrow{AC}$ in kot $\angle BAC=\alpha$, 

    %     \begin{solution}[5cm]

    % \end{solution}




    % \end{parts}







    % \newpage
    % 5. vprašanje

    \question
    Daljica $AB$ je določena s krajiščema $A(3,2,-2)$ in $B(1,1,-5)$.
    
    \begin{parts}
        \part[3]
        Določite koordinate točke $T$, ki leži na daljici $AB$, da bo veljalo $|AT|:|TB|=1:2$.
       
    %        \begin{solution}[5cm]

    % \end{solution}

        \part[3]    
        Določite enotski vektor v smeri krajevnega vektorja točke $B$.

    % \begin{solution}[5cm]

    % \end{solution}

        \part[3]
        Ali sta krajevna vektorja $\overrightarrow{r_A}$ in $\overrightarrow{r_B}$ pravokotna?

    %         \begin{solution}[5cm]

    % \end{solution}

        \part[3]
        Določite $u$, da bo vektor $\overrightarrow{v}=(u,u+3,9)$ kolinearen z vektorjem $\overrightarrow{AB}$.
    \end{parts}

% \begin{solution}[4cm]

% \end{solution}


% \newpage

    % 6. vprašanje



    \question[3]
    Dolžina vektorja $\overrightarrow{a}$ je $5$, dolžina vektorja $\overrightarrow{b}$ je $8$, 
    kot med njima pa meri $60^\circ$. 
    % Izračunajte skalarni produkt vektorjev $\overrightarrow{a}$ in $\frac{1}{2}\overrightarrow{a}+\frac{4}{3}\overrightarrow{b}$.
    % \begin{parts}
    %     \part[4]
        Izračunajte dolžino vektorja $\overrightarrow{c}=\frac{1}{5}\overrightarrow{a}+\frac{3}{2}\overrightarrow{b}$.
        
    %         \begin{solution}[5cm]

    % \end{solution}

    %     \part[3]
    %     Izračunajte skalarni produkt vektorjev $\overrightarrow{a}$ in $\frac{1}{2}\overrightarrow{a}+\frac{4}{3}\overrightarrow{b}$.

    % \end{parts}


    % \begin{solution}[10cm]

    % \end{solution}





\end{questions}



\end{document}